                                              IMC 2024
                                    First Day, August 7, 2024
                                                 Solutions


Problem 1. Determine all pairs (a, b) ∈ C × C satisfying

                                     |a| = |b| = 1 and a + b + ab ∈ R.

                                                            (proposed by Mike Daas, Universiteit Leiden)

Hint: Write a = eix and b = eiy , and transform the RHS to a product.

Solution 1. Write a = eix and b = eiy for some x, y ∈ [0, 2π). Using Euler’s formula, and the
well-known identities
                                           x+y     x−y                              x    x
                   sin x + sin y = 2 sin       cos           and    sin x = 2 sin     cos ,
                                            2       2                               2    2
we get a product form of the left-hand side as
                                             
                  Im a + b + ab = sin x + sin y + sin(x − y)
                                                x+y       x−y          x−y       x−y
                                        = 2 sin       cos      + 2 sin       cos
                                                 2         2           2        2
                                                  x+y         x−y        x−y
                                        = 2 sin         + sin        cos
                                                    2          2           2
                                                x      y      x−y
                                        = 4 sin · cos · cos       .
                                                2      2       2
Hence, a + b + ab is real if and only if either sin x2 = 0, cos y2 = 0 or cos x−y
                                                                               2
                                                                                  = 0, which respectivelly
correspond to x = 2kπ, y = (2k + 1)π and x = y + (2k + 1)π.
   Therefore, the solutions are

                         (1, b),    (a, −1) and (a, −a) with |a| = 1, |b| = 1.


Solution 2. Notice that

                                   a + b + ab ∈ R ⇐⇒ 1 + a + b + ab ∈ R.

Let c ∈ C be such that a = c2 . Now observe that

                                   c(1 + a + b + ab) = c + cc2 + cb + cc2 b
                                                     = c + c + cb + cb ∈ R,

where we used that cc = 1 and z + z ∈ R for any z ∈ C. W e conclude that either c ∈ R, or
1 + a + b + ab = 0. In the first case, c = ±1 and so a = 1. In the second case, we factor the equation
as
             (a + b)(1 + b) = 1 + a + 1b + ab = 0, and as such, a = −b or b = −1.
We find precisely three families of pairs (a, b): the pairs (1, b) for b on the unit circle; the pairs (a, −1)
for a on the unit circle; and the pairs (a, −a) for a on the unit circle.


                                                       1
Problem 2. For n = 1, 2, . . . let
                                             √
                                              2
                                                                            √
                                             n
                              Sn = log            11 · 22 · . . . · nn − log( n),

where log denotes the natural logarithm. Find lim Sn .
                                                      n→∞

                                     (proposed by Sergey Chernov, Belarusian State University, Minsk)

Hint: Sn is (close to) a Riemann sum of a certain integral.

Solution. Transform Sn as
                                         n
                                  1 X              1
                             Sn =        k log k −   log n
                                  n2 k=1           2
                                     n                    !
                                  1 X      k      k            1
                                =              log + log n    − log n
                                  n k=1 n         n            2
                                        n               n
                                     1Xk        k log n X    1
                                =            log + 2      k − log n
                                     n k=1 n    n   n k=1    2
                                        n
                                  1Xk        k log n
                                =         log +      .
                                  n k=1 n    n  2n


                    log n                              n k
                                                     1 P       k
Here the last term        converges to 0. The sum           log is a Riemann sum for the integrable
                     2n                             n k=1 n    n                      
                                                                       1 2      n−1
function f (x) = x log x on the segment [0, 1] with the uniform grid    , ,...,     , 1 . Therefore
                                                                       n n       n
                 n                 n      Z 1              2            1
              1Xk        k      1X       k                   x          x2      1
          lim         log = lim       f    =    x log xdx =     log x −       =− .
              n k=1 n    n      n k=1    n   0                2         4 0     4

                                    1
Hence, lim Sn exists, and lim Sn = − .
                                    4




                                                         2
Problem 3. For which positive integers n does there exist an n × n matrix A whose entries are all in
{0, 1}, such that A2 is the matrix of all ones?
                                (proposed by Alex Avdiushenko, Neapolis University Paphos, Cyprus)

Hint: Let J be the n × n matrix with all ones. Consider A3 = AJ = JA.

Solution. Answer: Such a matrix A exists if and only if n is a complete square.
   Let Jn be the n × n matrix with all ones, so A2 = Jn . Consider the equality

                                               A3 = AJn = Jn A.

In the matrix AJn , all columns are equal to the sum of colums in A, that is, the (i, j)th entry in AJn
is the number of ones in the ith row of A. Similarly, the (i, j)th entry in Jn A is the number of ones
in the jth column of A. These numbers must be equal, so A contains the same number of ones in
every row and every column. Let this common number be k; then AJn = Jn A = kJn .
    Now from
                       nJn = Jn2 = (A2 )2 = A(AJn ) = A(kJn ) = k(AJn ) = k 2 Jn
we can read n = k 2 , so n must be a complete square.
   It remains to show an example for a matrix A of order n = k 2 . For l = 0, 1, . . . , k − 1, let Bl
be the k × k matrix whose (i, j)th entry is 1 if j − i ≡ l (mod k) and 0 otherwise, i.e., Bl can be
obtained from the identity matrix by cyclically shifting the colums l times, and let
                                                                
                                       B0 B1 B2 . . . Bk−1
                                     B0 B1 B2 . . . Bk−1 
                                A =  ..                      ..  ;
                                                                
                                              ..   .. . .
                                      .       .    .     .    . 
                                       B0 B1 B2 . . . Bk−1

The (i, j)th block in A2 is
                                         
                                     Bj−1
                              Bk−1  ...  = (B0 + B1 + . . . + Bk−1 )Bj−1 = Jk Bj−1 = Jk ,
                                  
            B0 B1 . . .
                                          
                                        Bj−1

so this matrix indeed satisfies A2 = Jk2 .




                                                      3
Problem 4. Let g and h be two distinct elements of a group G, and let n be a positive integer.
Consider a sequence w = (w1 , w2 , . . .) which is not eventually periodic and where each wi is either g
or h. Denote by H the subgroup of G generated by all elements of the form wk wk+1 . . . wk+n−1 with
k ≥ 1. Prove that H does not depend on the choice of the sequence w (but may depend on n).
                                                     (proposed by Ivan Mitrofanov, Saarland University)

Solution. Let Xm denote the subset of G of products of the form g1 . . . gm , where each gi is either g
or h.
Lemma. For all j = 1, 2, . . . , n and for all a, b ∈ Xj the ratio a−1 b is contained in H.
Proof. Induction in j.
    We start with the base case j = 1. By the pigeonhole principle, there exist k < ℓ for which the
sequences (wk+1 , . . . , wk+n−1 ) and (wℓ+1 , . . . , wℓ+n−1 ) coincide. If wk+m = wℓ+m for all positive integer
m, then the sequence w is eventually periodic with period ℓ − k. Thus, there exists m > 0 for which
wk+m ̸= wℓ+m . We have m ⩾ n, so wk+m−i = wℓ+m−i for i = 1, 2, . . . , n − 1. Therefore, since the
products x = wk+m−n+1 . . . wk+m and y = wℓ+m−n+1 . . . wℓ+m both are elements of H, the subgroup
H contains their ratios x−1 y and y −1 x. These ratios are equal to g −1 h and h−1 g (in some order), that
finishes the proof for j = 1.
    Induction step from j − 1 to j, 2 ⩽ j ⩽ n. We say that an element a ∈ Xj is a g-element,
correspondingly an h-element, if it can be represented as a = ga1 , correspondingly a = ha1 , where
a1 ∈ Xj−1 . The ratio of two g-elements, or of two h-elements, is a ratio of two elements of Xj−1 , thus,
it is in H by the induction hypothesis. Since the property a−1 b ∈ H is an equivalence relation on
pairs (a, b), it suffices to find a g-element and h-element whose ratio is in H.
    Define k, ℓ, m, as in the base case. The subgroup H contains the products

                               v = wk+m−n+j . . . wk+m wk+m+1 . . . wk+m+j−1 ,
                               u = wℓ+m−n+j . . . wℓ+m wℓ+m+1 . . . wℓ+m+j−1 .

Their ratio u−1 v is a ratio of g-element and an h-element in Xj , since {wk+m , wℓ+m } = {g, h} and
wk+m−i = wℓ+m−i for all i = 1, 2, . . . , n − j.
   The Lemma for j = n yields that H is the subgroup of G generated by Xn , and this description
does not depend on w.




                                                       4
Problem 5. Let n > d be positive integers. Choose n independent, uniformly distributed random
points x1 , . . . , xn in the unit ball B ⊂ Rd centered at the origin. For a point p ∈ B denote by f (p)
the probability that the convex hull of x1 , . . . , xn contains p. Prove that if p, q ∈ B and the distance
of p from the origin is smaller than the distance of q from the origin, then f (p) ≥ f (q).
                                             (proposed by Fedor Petrov, St Petersburg State University)

Solution. By radial symmetry of the distribution, f (p) depends only on |op| (the distance between
o and p), so, we may assume that p lies on the segment between o and q. For points x1 , . . . , xn
and x ∈ B denote by fx (x1 , . . . , xn ) the indicator function of the event “x is in the convex hull of
x1 , . . . , xn ”. The claim follows from the following deterministic inequality
                                  X                            X
                                     fp (±x1 , . . . , ±xn ) ⩾   fq (±x1 , . . . , ±xn ),             (1)

where x1 , . . . , xn ∈ B are arbitrary points in general position and the summations are over all 2n choices
of signs (here o is identified with the origin, that is, x and −x are symmetric with respect to o). Indeed,
taking the expectation in (1) over independent random uniform x1 , . . . , xn , we get 2n f (p) ⩾ 2n f (q).
(To be specific, here “general position” means that for any point set A ⊂ {±x1 , . . . , ±xn , p, q}, which
does not contain simultaneosuly xi and −xi , is not contained in an (affine) (|A| − 2)-dimensional
plane. This holds with probability 1.)
    To prove (1), we use the following formula for the characteristic function χP of the convex
polyhedron
       P        P ⊂ Rd : if P1 , . . . , Pk are all facets of P , and Qi is the convex hull of o and Pi , then
χP =      ±χQi , where the sign is plus if o and P are on the same side of Pi , and minus otherwise.
Indeed, for every point p in general position look how the ray op intersects the boundary of P and
realize that for at most two summands the contribution of the RHS at point p is non-zero, and the
total contribution equals 1 when p is inside P and 0 (possibly as 0 = 1 − 1) otherwise. Use this
formula for every polyhedron P with n vertices y1 , . . . , yn , where each yi is ±xi . These polyhedrons
are simplicial (all facets are simplices) P    because of the general position condition. Sum up over all 2n
such P , we get the expression of P χP as a linear combination of χS , where S are simplices formed
by o and some d points in {±x1 , . . . , ±xn } (not containing xi and −xi simultaneously).
    For proving (1), it suffices to verify that all coefficients
                                                           P        of χS in this linear combination are positive
(since two sides of (1) are the values of the sum P χP at p and q). Let’s find a coefficient of χS ,
where, say, S is a simplex with vertices o, x1 , . . . , xd . The plane α through x1 , . . . , xd partitions Rd
onto two parts H + (containing o) and H − (not containing o). For every pair {xi , −xi } with i > d,
either both points belong to H + , or one belongs to H − and another to H + . χS goes with the plus
sign for P with vertices x1 , . . . , xd and other vertices from H + , and with the minus sign for P with
vertices x1 , . . . , xd and other vertices from H − . It is immediate that there are at least as many pluses
as minuses.




                                                       5
